\documentclass[journal]{IEEEtran}
\usepackage{graphicx}
\usepackage{amsmath}
\usepackage{booktabs}
\usepackage{cite}
\usepackage{url}
\usepackage{amssymb}

\begin{document}

\title{Multitask Spatiotemporal Joint Representation Learning for Echocardiography View Classification and Ejection Fraction Prediction}

\author{Amir~Khan\\
\normalsize Clinical Machine Learning Division, Deep Health Research Initiative\\
\normalsize Email: amir701771@users.noreply.github.com}

\maketitle

\begin{abstract}
Accurate assessment of Left Ventricular Ejection Fraction (LVEF) is vital for diagnosing cardiovascular diseases, yet manual evaluation suffers from significant inter-observer variability and high cognitive load. Furthermore, LVEF estimation is view-dependent, strictly requiring an Apical 4-Chamber (A4C) perspective. In this work, we propose a multitask deep learning framework that integrates view classification and LVEF regression. Using a shared spatiotemporal backbone (R2+1D) and a semi-supervised pseudo-labeling strategy, the model learns shared features that generalize across medical datasets (CAMUS, EchoNet-Dynamic, HMC-QU). Our model achieves a view classification accuracy of 98.2\% and an LVEF estimation Mean Absolute Error (MAE) of 4.12\%, demonstrating that joint multitask representation learning provides regularization that improves regression performance over single-task models.
\end{abstract}

\begin{IEEEkeywords}
Echocardiography, Multitask Learning, Ejection Fraction, View Classification, Deep Learning, Spatiotemporal representations.
\end{IEEEkeywords}

\section{Introduction}
\IEEEPARstart{L}{eft} ventricular ejection fraction (LVEF) is the primary clinical parameter utilized to assess cardiac systolic function. It measures the percentage of blood pumped out of the left ventricle during each contraction cycle. Quantitative estimation of LVEF is crucial for diagnostic classification, clinical staging, and management of congestive heart failure and myocardial infarction. However, standard clinical practices relying on visual estimation or manual boundary tracing (Biplane Simpson's method) are highly subjective, time-consuming, and subject to significant intra- and inter-observer variations (often up to 10\% absolute variance).

An inherent challenge in automated clinical diagnostic systems is view validation. Standard medical protocols prescribe distinct viewing windows---most notably the Apical 2-Chamber (A2C) and Apical 4-Chamber (A4C) orientations---to assess localized wall motion anomalies. LVEF calculations are highly sensitive to these orientations, and calculating ejection fraction on misaligned or incorrect view planes leads to severe diagnostic error. Therefore, a complete clinical tool must first validate the viewing angle before performing regression. Integrating these tasks in a unified multitask framework provides a robust solution, saving computational overhead and regularizing feature representations.

\section{Related Work}
Prior work in automated echocardiography has focused predominantly on single-task architectures. View classification systems typically employ 2D Convolutional Neural Networks (CNNs) trained on isolated frames, ignoring key spatiotemporal cues present in cardiac contraction cycles. Conversely, LVEF estimation models, such as EchoNet-Dynamic, leverage 3D CNNs trained directly on A4C video series. However, these regression networks ignore the broader spatial context and lack the capability to validate input orientation automatically. Recent studies in multitask learning suggest that training deep networks on complementary auxiliary tasks, such as predicting view boundaries alongside numerical regression, forces the shared backbone layers to build more generalized feature maps, avoiding overfitting to noise.

\section{Methodology}

\subsection{Dataset Preprocessing}
Raw echocardiograms are stored in multi-frame video formats (AVI/DICOM) or medical imaging formats (MHD/NIFTI). The preprocessing pipeline uniformly samples a fixed sequence of frames ($T=16$) to capture the complete cardiac cycle. Frames are resized to $112 \times 112$ pixels and normalized using standard ImageNet mean and standard deviation values. Data augmentation, including random rotations, translations, and horizontal flips, is applied during training to prevent overfitting.

\subsection{Model Architecture}
The core model, MultiTaskCNN, consists of a shared spatiotemporal backbone followed by two independent head branches. We compare 2D backbones (ResNet-18) using frame-averaged pooling against 3D convolutional backbones ($R(2+1)D$-18). The shared feature extractor outputs a 512-dimensional representation vector, which is passed through a Dropout layer ($p=0.3$ or $0.4$). The Classification Head maps the features to 2 classes (A2C and A4C views) via a linear layer. The Regression Head uses a linear regressor with a clinical clamping function that bounds predicted LVEF values within the physiological limits:
\begin{equation}
\text{LVEF}_{\text{clamped}} = \max(10.0, \min(95.0, \text{LVEF}_{\text{raw}}))
\end{equation}

\subsection{Loss Formulation \& Hyperparameter Tuning}
The model is optimized using a joint loss function:
\begin{equation}
\mathcal{L}_{\text{total}} = \alpha \mathcal{L}_{\text{view}} + \beta \mathcal{L}_{\text{ef}}
\end{equation}
where $\mathcal{L}_{\text{view}}$ is Weighted Binary Cross-Entropy (compensating for class imbalance) and $\mathcal{L}_{\text{ef}}$ is Mean Squared Error (MSE). We utilize an evolutionary algorithm (DEAP) to automatically optimize hyperparameters. The chromosome encodes the loss weights $\alpha$ and $\beta$, the initial learning rate, and weight decay. The fitness function is evaluated on a validation set, targeting maximum classification accuracy and minimal regression MAE.

\section{Experiments \& Results}

\subsection{Experimental Setup}
All models were implemented in PyTorch and trained on NVIDIA RTX GPUs. Pre-training of the view classifier was performed on the CAMUS dataset. Semi-supervised pseudo-labeling was executed on the unlabelled frames of the EchoNet-Dynamic dataset, enabling joint multitask training on 10,036 patients. Models were trained for 50 epochs using Adam optimizer, with a batch size of 4 and learning rate initialized at $0.0001$.

\subsection{Model Comparison}
The quantitative results are summarized in Table I. The joint multitask model consistently outperforms the single-task variants. In particular, joint training with $R(2+1)D$-18 achieves a view classification accuracy of 98.2\% and an LVEF prediction MAE of 4.12\%.

\begin{table}[htbp]
\caption{Performance Benchmarking on the Validation Sets}
\label{table_results}
\centering
\begin{tabular}{llcccc}
\toprule
\textbf{Model Configuration} & \textbf{Backbone} & \textbf{View Acc (\%)} & \textbf{LVEF MAE (\%)} & \textbf{LVEF RMSE (\%)} & \textbf{LVEF $R^2$} \\
\midrule
Single-Task (View Only) & ResNet-18 & 95.8\% & -- & -- & -- \\
Single-Task (LVEF Only) & ResNet-18 & -- & 5.34\% & 6.82\% & 0.68 \\
\textbf{Multi-Task (Joint)} & \textbf{ResNet-18} & \textbf{97.1\%} & \textbf{4.65\%} & \textbf{5.98\%} & \textbf{0.75} \\
\midrule
Single-Task (View Only) & R2+1D & 96.9\% & -- & -- & -- \\
Single-Task (LVEF Only) & R2+1D & -- & 4.88\% & 6.12\% & 0.74 \\
\textbf{Multi-Task (Joint)} & \textbf{R2+1D} & \textbf{98.2\%} & \textbf{4.12\%} & \textbf{5.11\%} & \textbf{0.81} \\
\bottomrule
\end{tabular}
\end{table}

\section{Discussion \& Conclusion}
Our findings confirm that multitask representation learning acts as an effective regularizer in cardiac ultrasound analysis. The inclusion of view classification forces the feature extractor to capture structural patterns like chamber alignment and valve movement, which directly benefits the LVEF regression head. Furthermore, spatiotemporal backbones show a substantial advantage over frame-aggregated 2D CNNs, reducing LVEF MAE by 11.4\%. Future directions include extending this multitask framework to segment the left ventricular wall and calculate longitudinal strain parameters, paving the way for fully automated point-of-care cardiac analysis.

\begin{thebibliography}{99}
\bibitem{ouyang2020} D. Ouyang et al., ``Video-based AI for beat-to-beat assessment of cardiac function,'' \emph{Nature}, vol. 580, pp. 252--256, 2020.
\bibitem{leclerc2019} S. Leclerc et al., ``Deep learning for segmentation using the CAMUS dataset,'' \emph{IEEE Trans. Med. Imaging}, vol. 38, pp. 2198--2210, 2019.
\bibitem{demir2021} A. Demir et al., ``HMC-QU Dataset for myocardial infarction detection,'' \emph{IEEE Access}, vol. 9, pp. 12431--12440, 2021.
\bibitem{tran2018} D. Tran et al., ``A closer look at spatiotemporal convolutions for video recognition,'' \emph{CVPR}, pp. 6450--6459, 2018.
\end{thebibliography}

\end{document}
